
\documentclass{article}
\usepackage{colortbl}
\usepackage{makecell}
\usepackage{multirow}
\usepackage{supertabular}

\begin{document}

\newcounter{utterance}

\centering \large Interaction Transcript for game `codenames', experiment `risk\_high', episode 8 with reflexion{-}Qwen3.5{-}27B.
\vspace{24pt}

{ \footnotesize  \setcounter{utterance}{1}
\setlength{\tabcolsep}{0pt}
\begin{supertabular}{c@{$\;$}|p{.15\linewidth}@{}p{.15\linewidth}p{.15\linewidth}p{.15\linewidth}p{.15\linewidth}p{.15\linewidth}}
    \# & $\;$A & \multicolumn{4}{c}{Game Master} & $\;\:$B\\
    \hline

    \theutterance \stepcounter{utterance}  
    & & & \multicolumn{2}{p{0.3\linewidth}}{
        \cellcolor[rgb]{0.9,0.9,0.9}{
            \makecell[{{p{\linewidth}}}]{
                \texttt{\tiny{[GM$|$GM]}}
                {'hidden': {'team': ['octopus'], 'innocent': ['hook', 'robot'], 'opponent': [], 'assassin': ['moscow', 'tick', 'forest', 'circle']}, 'revealed': {'team': {'team': [], 'innocent': [], 'opponent': [], 'assassin': []}, 'opponent': {'team': [], 'innocent': [], 'opponent': [], 'assassin': []}}}
            }
        }
    }
    & & \\ \\

    \theutterance \stepcounter{utterance}  
    & & \multicolumn{4}{p{0.6\linewidth}}{
        \cellcolor[rgb]{0.9,0.9,0.9}{
            \makecell[{{p{\linewidth}}}]{
                \texttt{\tiny{[P1$\langle$GM]}}
                \texttt{You have attempted this same task before and did not succeed. These are the plans you wrote after those attempts.} \\
\texttt{Trial \#1: New plan: My previous strategy failed because I prioritized covering multiple words with a single clue over ensuring the clue was unambiguous, leading my partner to guess an opponent's word ("temple" for "mount" and likely "mine" for "sour" due to the "acid mine" association). Next time, I will prioritize safety and clarity over speed. I will select clues that have a very strong, exclusive semantic link to my target words and a weak or non{-}existent link to any opponent, distractor, or kill words. Specifically, I will avoid clues like "sour" that could easily trigger "mine" (acid mine) or "forest" (sour forest?), and instead use more direct associations like "citrus" for "lemon" or "medical" for "doctor," even if it means guessing fewer words per turn. I will also explicitly verify that my chosen clue does not share a common compound word or idiom with any of the opponent's remaining words before submitting.} \\
\\ 
\texttt{Trial \#2: New plan: My previous strategy failed because I prioritized speed by grouping words (e.g., "air" and "fair" with "WIND") without sufficiently verifying that the clue wouldn't trigger a distractor word. The clue "WIND" successfully linked to my targets but also strongly associated with the distractor "bugle" (as in "wind instrument"), causing my turn to end prematurely. Next time, I will prioritize safety over grouping. Before submitting a clue, I will explicitly cross{-}reference the clue against all remaining distractor, opponent, and kill words to ensure no semantic overlap exists. If a clue risks triggering a distractor, I will split the targets into separate turns or choose a more specific, exclusive clue, even if it means guessing fewer words per turn. For example, instead of "WIND" for "air" and "fair", I would have used "BREEZE" for "air" and "JUST" for "fair" in separate turns to avoid the "bugle" association.} \\
\\ 
\texttt{Trial \#3: New plan: My previous strategy failed because I chose a clue ("TENTACLE") that, while semantically related to my target "octopus," was too specific and potentially ambiguous in the context of the distractor "robot" (which can have tentacle{-}like appendages in sci{-}fi contexts) or simply failed to provide a strong enough unique association to prevent a misinterpretation. Although the game ended in a loss, the primary error was likely the lack of a broader, more distinct semantic category that clearly separates "octopus" from "robot" and "hook." Next time, I will select a clue that relies on a fundamental biological or categorical distinction that "robot" and "hook" cannot share. For "octopus," I will use a clue like "CEPHALOPOD" or "INK" or "SEA" (if safe from other sea{-}related distractors) to ensure the association is exclusive to the animal kingdom and distinct from mechanical or tool{-}based distractors. I will explicitly verify that the clue does not trigger any association with "robot" (mechanical) or "hook" (tool/fishing) before submitting. If a single word cannot safely cover the target without risk, I will prioritize the safest, most exclusive association over speed.} \\
\\ 
\\ 
\texttt{Provide one single{-}word clue that relates to some of your team's words from the list below. You can choose to how many and to which words your clue relates to.} \\
\texttt{The goal is to have your partner guess all of your team's words based on the clues you give them, before the other team has uncovered all of theirs. If your partner guesses an 'opponent' or a 'distractor' word, their guessing turn ends and the opposing team gets to make a turn. If your partner guesses a 'kill' word, the game ends and your team loses immediately.} \\
\texttt{The more words your clue relates to, the quicker you can win the game, but the harder the associations may be to guess for your partner, so choose your clue word wisely. The clue word has to be semantically related to the target words, it cannot be one of the words in the lists or contain parts of them.} \\
\\ 
\texttt{Always give your single{-}word clue and your comma{-}separated list of related target words in the following format and make your answers as short as possible, never include any other text than is required in this form:} \\
\\ 
\texttt{CLUE: <WORD>} \\
\texttt{TARGETS: <TARGETS>} \\
\\ 
\texttt{Your team words are: octopus.} \\
\texttt{Your opponent's team words are: .} \\
\texttt{Distractor words are: hook, robot.} \\
\texttt{The kill words are: moscow, tick, forest, circle.} \\
            }
        }
    }
    & \\ \\

    \theutterance \stepcounter{utterance}  
    & \multicolumn{4}{p{0.6\linewidth}}{
        \cellcolor[rgb]{0.9,0.9,0.9}{
            \makecell[{{p{\linewidth}}}]{
                \texttt{\tiny{[P1$\rangle$GM]}}
                \texttt{CLUE: INK} \\
\texttt{TARGETS: octopus} \\
            }
        }
    }
    & & \\ \\

    \theutterance \stepcounter{utterance}  
    & & & \multicolumn{2}{p{0.3\linewidth}}{
        \cellcolor[rgb]{0.9,0.9,0.9}{
            \makecell[{{p{\linewidth}}}]{
                \texttt{\tiny{[GM$|$GM]}}
                \texttt{ink} \\
            }
        }
    }
    & & \\ \\

    \theutterance \stepcounter{utterance}  
    & & & \multicolumn{2}{p{0.3\linewidth}}{
        \cellcolor[rgb]{0.9,0.9,0.9}{
            \makecell[{{p{\linewidth}}}]{
                \texttt{\tiny{[GM$|$GM]}}
                ['octopus']
            }
        }
    }
    & & \\ \\

    \theutterance \stepcounter{utterance}  
    & & & \multicolumn{2}{p{0.3\linewidth}}{
        \cellcolor[rgb]{0.9,0.9,0.9}{
            \makecell[{{p{\linewidth}}}]{
                \texttt{\tiny{[GM$|$GM]}}
                {'word': 'octopus', 'assignment': 'team'}
            }
        }
    }
    & & \\ \\

    \theutterance \stepcounter{utterance}  
    & & & \multicolumn{2}{p{0.3\linewidth}}{
        \cellcolor[rgb]{0.9,0.9,0.9}{
            \makecell[{{p{\linewidth}}}]{
                \texttt{\tiny{[GM$|$GM]}}
                \texttt{opponent has won} \\
            }
        }
    }
    & & \\ \\

    \theutterance \stepcounter{utterance}  
    & \multicolumn{4}{p{0.6\linewidth}}{
        \cellcolor[rgb]{0.95,0.95,0.95}{
            \makecell[{{p{\linewidth}}}]{
                \texttt{\tiny{[P2$\rightarrow$P1]}}
                \texttt{Post{-}episode Reflexion (not sent to GM):} \\
\texttt{New plan: My previous attempt failed because the clue "INK" was too broad and ambiguous; while octopuses produce ink, the word "ink" is also strongly associated with writing, printing, and potentially the "robot" distractor (e.g., inkjet printers or robotic arms in manufacturing), leading to a misinterpretation or a failed guess. I failed to verify that the clue was exclusive to the biological target and distinct from the mechanical distractor. Next time, I will select a clue that relies on a fundamental biological or habitat{-}based distinction that "robot" and "hook" cannot share. For "octopus," I will use a clue like "CEPHALOPOD" (a specific biological class) or "EIGHT" (referring to its eight arms, distinct from a robot's variable limbs or a hook's single curve) to ensure the association is exclusive. I will explicitly verify that the chosen clue has no semantic overlap with "robot" (mechanical) or "hook" (tool) before submitting, prioritizing a unique, unambiguous link over a common but risky association.} \\
            }
        }
    }
    & & \\ \\

\end{supertabular}
}

\end{document}
